\documentclass[11pt,a4paper]{article}
\usepackage[T1]{fontenc}
\usepackage{lmodern}
\usepackage{amsmath,amssymb,amsthm}
\usepackage[margin=28mm]{geometry}
\usepackage{microtype}
\usepackage[hidelinks]{hyperref}
\hypersetup{pdftitle={Closed physical realizations and a compact-resolvent pencil for a theta prime-Gamma energy},pdfauthor={Jongmin Choi},pdfsubject={Conditional physical realization and exact domain analysis},pdfkeywords={theta kernel, closed quadratic form, compact resolvent, prime powers, spectral pencil}}
\setlength{\parindent}{1.2em}
\setlength{\parskip}{0.25em}
\setlength{\emergencystretch}{2em}
\allowdisplaybreaks[2]
\numberwithin{equation}{section}
\newtheorem{theorem}{Theorem}[section]
\newtheorem{proposition}[theorem]{Proposition}
\newtheorem{lemma}[theorem]{Lemma}
\newtheorem{corollary}[theorem]{Corollary}
\theoremstyle{definition}
\newtheorem{hypothesis}[theorem]{Hypothesis}
\theoremstyle{remark}
\newtheorem{remark}[theorem]{Remark}
\newcommand{\R}{\mathbb R}
\newcommand{\C}{\mathbb C}
\newcommand{\E}{\mathcal E}
\newcommand{\D}{\mathcal D}
\newcommand{\He}{L^2_{\mathrm{even}}(\mu)}
\newcommand{\Le}{L^2_{\mathrm{even}}(dx)}
\newcommand{\tg}{t_{\Gamma}}
\newcommand{\cg}{c_{\Gamma}}
\newcommand{\dom}{\operatorname{Dom}}
\newcommand{\Var}{\operatorname{Var}}
\newcommand{\rank}{\operatorname{rank}}
\newcommand{\ess}{\operatorname{ess}}
\newcommand{\supp}{\operatorname{supp}}
\newcommand{\ip}[2]{\langle #1,#2\rangle}
\newcommand{\norm}[1]{\lVert #1\rVert}
\newcommand{\one}{\mathbf 1}
\title{Closed physical realizations and a compact-resolvent pencil for a theta prime--Gamma energy}
\author{Jongmin Choi\\[3pt]\small Independent Researcher, Seoul, Korea\\\small ORCID iD: \href{https://orcid.org/0009-0008-7448-514X}{0009-0008-7448-514X}}
\date{November 1, 2026\\[3pt]\small Draft v0.1}
\begin{document}
\maketitle
\begin{abstract}
We construct a closed self-adjoint physical pencil for a theta-weighted energy containing the complete prime-power sum. A double-exponential theta envelope gives an operator-norm estimate for the weighted prime tail. Together with logarithmic Fourier coercivity of the Gamma form, it yields a common form domain and compact resolvent for each positive pencil parameter. Under the explicitly stated compact physical bridge hypothesis \textup{(H)}, multiplication by the theta kernel identifies this domain exactly with the domain of the original closed energy. The resulting kernel correspondence concerns actual operator-domain eigenstates, and a form congruence gives the precise negative index. We also prove that the original mean is unbounded in ordinary physical $L^2$, and that the theta kernel is not in the pencil operator domain except at the distinguished parameter $1/2$. These facts prevent an ordinary-$L^2$ ground-vector removal argument. The unresolved global variance bound is equivalent, under \textup{(H)}, to strict positivity of the pencil's second eigenvalue at every parameter in $(0,1/2)$. That sign is not proved. The paper neither proves the global half bound nor the Riemann hypothesis.
\end{abstract}

\section{Introduction and principal result}

A change of dependent variable can make an energy operator appear confining without preserving its Hilbert space, operator domain, or spectral decomposition. This issue is particularly consequential for theta-weighted prime--Gamma energies. Their positive integral form is naturally defined in a rapidly decaying probability measure, whereas the physical numerator is often written in ordinary Lebesgue $L^2$. The two realizations must be connected through closed forms, rather than by treating their formal rows as interchangeable operators.

This paper establishes that connection under a precisely isolated compact-test identity. The physical pencil itself is constructed without assuming the identity. Its compact resolvent follows from two distinct mechanisms: the theta weight controls spatial escape, and the Gamma difference form controls high Fourier frequencies. The complete prime-power sum is retained throughout. Finite prime heads enter only in estimates with a proved infinite remainder.

Write $c(x)=\cosh(x/2)$ and let $\kappa$ be the positive even theta kernel defined in Section~\ref{sec:energy}. Put
\begin{equation}\label{eq:weights}
 d\mu(x)=2c(x)\kappa(x)\,dx,\qquad
 W(x)=\frac{\kappa(x)}{2c(x)},\qquad
 \norm{u}_\omega^2=\int_\R\frac{|u(x)|^2}{W(x)}\,dx.
\end{equation}
The measure $\mu$ has mass one. The full positive energy $\E$, its closed domain $\D$, and its self-adjoint operator $A$ will be defined directly from nonnegative integrals. The centered restriction is denoted by $A_E$.

The physical Gamma form and the complete prime correlation form are
\begin{align}
 \tg[u]&=\int_0^\infty\frac{e^{-s/2}}{1-e^{-2s}}
                \norm{u(\cdot+s)-u}_2^2\,ds,\label{eq:physicalgamma}\\
 p[u]&=2\sum_{n\ge2}\frac{\Lambda(n)}{\sqrt n}
             \operatorname{Re}\int_\R u(x+s_n)\overline{u(x)}\,dx,
       \qquad s_n=\log n.\label{eq:physicalprime}
\end{align}
Here $\Lambda$ is the von Mangoldt function: $\Lambda(p^k)=\log p$ for primes $p$ and integers $k\ge1$, and $\Lambda(n)=0$ otherwise. Set
\begin{equation}\label{eq:qzero}
 \cg=\log\pi-\psi(1/4),\qquad
 q_0[u]=\tg[u]-\cg\norm{u}_2^2-p[u],
\end{equation}
where $\psi=\Gamma'/\Gamma$. The correlation series is initially defined on even compact smooth functions; Section~\ref{sec:prime} extends it absolutely to every $u$ of finite weighted norm.

\begin{hypothesis}[Compact physical bridge {\rm(H)}]\label{hyp:bridge}
For every even $f\in C_c^\infty(\R)$,
\begin{equation}\tag{H}\label{eq:bridge}
 q_0[\kappa f]=\E[f]-\frac12\norm{f}_\mu^2.
\end{equation}
\end{hypothesis}

Hypothesis~\ref{hyp:bridge} is an assumption of the energy-identification and spectral-correspondence results in this paper. It is not derived here from an explicit formula, a theta null identity, or a zero-location assertion. In particular, the normalization of $\mu$ and the theta envelope do not by themselves establish~\eqref{eq:bridge}. All uses of this hypothesis are identified below.

\begin{theorem}[Closed physical realization and spectral correspondence]\label{thm:main}
Let $\kappa$, $\mu$, $W$, and $\E$ be as defined above and in Section~\ref{sec:energy}. Define
\begin{equation}\label{eq:X}
 X=\{u\in\Le:\tg[u]<\infty,\ \norm{u}_\omega<\infty\},\qquad
 \norm{u}_X^2=\tg[u]+\norm{u}_\omega^2.
\end{equation}
For every $\delta>0$, the form
\begin{equation}\label{eq:ell}
 \ell_\delta[u]=q_0[u]+\delta\norm{u}_\omega^2,
 \qquad\dom(\ell_\delta)=X,
\end{equation}
is densely defined, closed, and lower bounded on $\Le$. Its self-adjoint representative $L_\delta$ has compact resolvent. These conclusions do not require \textup{(H)}.

Assume in addition \textup{(H)}. Then $f\mapsto\kappa f$ is a bijection $\D\to X$ and extends the bridge to the full form domain. For $0<\delta<1/2$ and $\lambda=1/2-\delta$,
\begin{align}
 \ker L_\delta
 &=\{\kappa f:f\in\dom(A_E),\ A_Ef=\lambda f\},\label{eq:mainkernel}\\
 n_-(L_\delta)
 &=1+\rank\one_{(0,\lambda)}(A_E).\label{eq:mainindex}
\end{align}
The centered operator $A_E$ has a strictly positive lower bound and every spectral projection $\one_{(-\infty,t]}(A_E)$, $t<1/2$, has finite rank. Moreover, if $\eta_1(\delta)\le\eta_2(\delta)\le\cdots$ are the eigenvalues of $L_\delta$ with multiplicity, then
\begin{equation}\label{eq:mainequivalence}
 \E[f]\ge\tfrac12\Var_\mu(f)\quad(f\in\D)
 \quad\Longleftrightarrow\quad
 \eta_2(\delta)>0\quad(0<\delta<1/2).
\end{equation}
Neither side of~\eqref{eq:mainequivalence} is established by the construction.
\end{theorem}

Here $n_-$ is the number of negative eigenvalues, counted with multiplicity; equivalently it is the maximal dimension of a strictly negative form subspace. Sections~\ref{sec:domain}--\ref{sec:sign} prove the theorem. Section~\ref{sec:mean} explains why the mean-zero condition cannot be replaced by ordinary-$L^2$ orthogonality to $\kappa$.

All Hilbert spaces are complex unless stated otherwise, and inner products are linear in the first argument. Quadratic forms are polarized accordingly. The real-even subspaces are invariant. No eigenvalue computation, prime enumeration, or numerical sign certificate is used. The paper makes no claim of historical priority.

\section{The theta kernel and the original closed energy}\label{sec:energy}

For $x\ge0$, define
\begin{equation}\label{eq:kappa}
 \kappa(x)=e^{x/2}\sum_{j\ge1}(4\tau_j(x)^2-6\tau_j(x))e^{-\tau_j(x)},
 \qquad \tau_j(x)=\pi j^2e^{2x},
\end{equation}
and extend evenly. We record the normalization and decay needed later.

\begin{lemma}[Smoothness and normalization]\label{lem:theta}
The kernel $\kappa$ is smooth, strictly positive, and even. There is a finite $C_\kappa>0$ such that
\begin{equation}\label{eq:envelopek}
 \kappa(x)\le C_\kappa\exp\!\left(\frac92|x|-\pi e^{2|x|}\right),
 \qquad
 C_\kappa=4\pi^2\sum_{j\ge1}j^4e^{-\pi(j^2-1)}.
\end{equation}
Furthermore $\int_\R2c\kappa\,dx=1$, and
\begin{equation}\label{eq:envelopeW}
 0<W(x)\le C_\kappa\exp(4|x|-\pi e^{2|x|}).
\end{equation}
\end{lemma}
\begin{proof}
The series in~\eqref{eq:kappa}, together with every derivative, converges uniformly on compact subsets of $[0,\infty)$. Each summand is positive because $\tau_j\ge\pi>3/2$. Dropping the negative polynomial term and using $e^{2x}\ge1$ gives~\eqref{eq:envelopek} for $x\ge0$.

To justify smooth evenness at zero, rather than assume it from reflection, set
\[
 H(x)=e^{x/2}\sum_{j\ge1}e^{-\pi j^2e^{2x}},\qquad x\in\R.
\]
The Jacobi transformation \cite[Eq.~20.7.32]{DLMF} implies
$H(-x)=H(x)+\sinh(x/2)$. Applying $D^2-1/4$ shows that $H''-H/4$ is even; direct differentiation identifies it with~\eqref{eq:kappa} on the positive half-line. Thus it gives the claimed smooth extension. Differentiating the transformation at zero yields $H'(0)=-1/4$. Since $c''=c/4$, integration by parts and the rapid decay of $H,H'$ at $+\infty$ give
\[
 \int_0^\infty c(x)\kappa(x)\,dx
  =[cH'-c'H]_0^\infty=\frac14.
\]
Evenness proves the normalization. Finally $2c(x)\ge e^{|x|/2}$ gives~\eqref{eq:envelopeW}.
\end{proof}

Put $\rho(s)=e^{-s/2}/(1-e^{-2s})$ for $s>0$, and $a_n=\Lambda(n)/\sqrt n$. The full positive energy is
\begin{align}
 \E_\Gamma[f]&=\int_0^\infty\rho(s)\int_\R
       \kappa(x)\kappa(x+s)|f(x+s)-f(x)|^2\,dx\,ds,\label{eq:Eg}\\
 \E_P[f]&=\sum_{n\ge2}a_n\int_\R
       \kappa(x)\kappa(x+s_n)|f(x+s_n)-f(x)|^2\,dx,\label{eq:Ep}\\
 \E[f]&=\E_\Gamma[f]+\E_P[f],\qquad
 \D=\{f\in\He:\E[f]<\infty\}.\label{eq:Edomain}
\end{align}
These are nonnegative integrals and sums, interpreted by Tonelli's theorem; no cancellation defines their domain. Write
\[
 \mu(f)=\int_\R f\,d\mu,\qquad
 \Var_\mu(f)=\norm{f}_\mu^2-|\mu(f)|^2.
\]

\begin{proposition}[Maximal domain and compact smooth form core]\label{prop:core}
The form $\E$ on $\He$ is densely defined, nonnegative, and closed. Its maximal integral domain $\D$ is the closure of even $C_c^\infty(\R)$ in the norm $(\norm{f}_\mu^2+\E[f])^{1/2}$. Constants belong to $\D$, and $\E[f]=0$ holds precisely for almost-everywhere constant $f$.
\end{proposition}
\begin{proof}
The theta envelope implies that for every $A_0>0$ there is a constant $C_{A_0}$ such that $\kappa(x)\le C_{A_0}e^{-A_0|x|}$. Consequently
\begin{equation}\label{eq:edgefinite}
 \int_\R\kappa(x)\kappa(x+s)\,dx
 \le C_{A_0}^2(s+A_0^{-1})e^{-A_0s},\qquad s\ge0.
\end{equation}
Since $a_n\le n$, the complete prime edge measure has finite mass, by taking $A_0>2$ in~\eqref{eq:edgefinite}. For every $L<\infty$,
\begin{equation}\label{eq:rhointegrable}
 \int_0^\infty\rho(s)\min\{1,L^2s^2\}\,ds<\infty.
\end{equation}
These observations admit compact smooth functions and constants to the domain.

For closedness, let $(f_j)$ be Cauchy in the form norm, with $L^2(\mu)$ limit $f$. A subsequence converges almost everywhere for $\mu$, hence for Lebesgue measure, because the density is strictly positive. A countable union of translated null sets is null, so all prime differences converge almost everywhere. Fubini's theorem gives convergence for the Gamma edge measure as well. For each fixed $k$, Fatou's lemma gives
\[
 \E[f-f_k]\le\liminf_{j\to\infty}\E[f_j-f_k].
\]
It follows that $f\in\D$ and $f_k\to f$ in form norm.

We prove the core assertion first for real profiles. Let $T_M(t)=\max\{-M,\min\{t,M\}\}$. The remainder $r_M(t)=t-T_M(t)$ is $1$-Lipschitz, tends pointwise to zero, and satisfies $|r_M(t)|\le|t|$. Thus dominated convergence in both edge measures and in $L^2(\mu)$ gives $T_M(f)\to f$ in form norm. It suffices to approximate bounded even profiles.

Choose smooth even cutoffs $0\le\chi_R\le1$ equal to one on $[-R,R]$, supported in $[-2R,2R]$, and with uniformly bounded Lipschitz constants. Equation~\eqref{eq:rhointegrable}, the finite prime-edge mass, and dominated convergence give $\E[\chi_R]\to0$. For bounded $f$, the identity for a product difference yields
\begin{align*}
 |[f(1-\chi_R)](y)-[f(1-\chi_R)](x)|^2
 &\le2(1-\chi_R(y))^2|f(y)-f(x)|^2\\
 &\quad+2\norm{f}_\infty^2|\chi_R(y)-\chi_R(x)|^2.
\end{align*}
Hence $f\chi_R\to f$ in form norm. It remains to smooth a compactly supported profile $g\in\D$.

Suppose $\supp g\subset[-K,K]$. Its ordinary $L^2$ norm is finite, since the density of $\mu$ is bounded below on this interval. For $0<s\le1$, a nonzero difference of $g$ has both endpoints in $[-K-1,K+1]$. The positive lower bound of $\kappa$ there controls the small-shift part of $\tg[g]$ by $\E_\Gamma[g]$. The part with $s>1$ is bounded by $4\norm{g}_2^2\int_1^\infty\rho(s)\,ds$. Thus $\tg[g]<\infty$.

Let $g_\epsilon=\zeta_\epsilon*g$, where $\zeta_\epsilon$ is an even compactly supported smooth approximate identity. Translations commute with convolution, and
\[
 \norm{\Delta_s(g_\epsilon-g)}_2\le2\norm{\Delta_sg}_2,
 \qquad\Delta_sg=g(\cdot+s)-g.
\]
Dominated convergence gives $\tg[g_\epsilon-g]\to0$. Since $\kappa$ is bounded,
$\E_\Gamma[g_\epsilon-g]\le\norm{\kappa}_\infty^2\tg[g_\epsilon-g]\to0$.
All the smoothed profiles have support in one fixed compact interval $J$. For any $h$ supported there,
\begin{equation}\label{eq:compactprimerate}
 \E_P[h]\le2\int_J\kappa(x)|h(x)|^2
     \sum_{n\ge2}a_n\bigl(\kappa(x-s_n)+\kappa(x+s_n)\bigr)\,dx
 \le C_J\norm{h}_2^2.
\end{equation}
Indeed, the complete rate in parentheses is uniformly bounded on $J$: use $\kappa(x\pm\log n)\le C_{A_0}e^{A_0\sup_J|x|}n^{-A_0}$ and $A_0>2$. Ordinary $L^2$ convergence therefore proves prime-energy convergence. The local upper bound of the density gives $L^2(\mu)$ convergence. Applying the argument to real and imaginary parts proves the complex core assertion.

Finally $\E[1]=0$. If $\E[f]=0$, the strictly positive Gamma density implies $f(y)=f(x)$ for almost every pair $x<y$. Fubini's theorem forces $f$ to be almost everywhere constant. Density of the core in $\He$ is immediate from the positive smooth density of $\mu$.
\end{proof}

The representation theorem for closed nonnegative forms \cite[Chapter~VI, Section~2]{Kato} defines $A$ by
\begin{equation}\label{eq:Adef}
 \begin{split}
 f\in\dom(A),\ Af=h\quad\Longleftrightarrow\quad
 f\in\D,\ h\in\He,\quad
 \E(f,g)=\ip{h}{g}_\mu\ \text{for all }g\in\D.
 \end{split}
\end{equation}
The operator is self-adjoint, $A1=0$, and $\ker A=\C1$. The centered space
$H_0=\{f\in\He:\mu(f)=0\}$ reduces $A$; its restriction is $A_E$. This construction precedes and does not use hypothesis~(H).

\section{The complete weighted prime form}\label{sec:prime}

The relevant prime bound is in the numerator's weighted norm, not a uniform vanishing estimate for an incoming jump rate. Let $\beta=\sqrt W$ and let $\tau_sh(x)=h(x+s)$ on ordinary $L^2(\R)$.

\begin{lemma}[Weighted shift estimate]\label{lem:prime}
For a finite constant $K$ and every $s\ge0$,
\begin{equation}\label{eq:betaprod}
 \sup_{x\in\R}\beta(x)\beta(x+s)\le K e^{2s-\pi e^s}.
\end{equation}
The self-adjoint series
\begin{equation}\label{eq:KP}
 K_P=\sum_{n\ge2}a_nM_\beta(\tau_{s_n}+\tau_{-s_n})M_\beta
\end{equation}
converges absolutely in operator norm and preserves even functions. If $K_{P,N}$ denotes its head through $n=N$, then
\begin{equation}\label{eq:rN}
 \norm{K_P-K_{P,N}}\le r_N,
 \qquad r_N=2K\sum_{n>N}\Lambda(n)n^{3/2}e^{-\pi n}\longrightarrow0.
\end{equation}
For every $u$ of finite weighted norm the series~\eqref{eq:physicalprime} converges absolutely and
\begin{align}
 p[u]&=\ip{K_P(u/\beta)}{u/\beta}_2,\label{eq:prepresentation}\\
 |p[u]-p_N[u]|&\le r_N\norm{u}_\omega^2,\qquad
 |p_N[u]|\le2\sum_{2\le n\le N}a_n\norm{u}_2^2.\label{eq:headtail}
\end{align}
\end{lemma}
\begin{proof}
Set $z=|x|+|x+s|\ge s$. Equation~\eqref{eq:envelopeW} and the arithmetic--geometric mean inequality give
\begin{align*}
 \beta(x)\beta(x+s)
 &\le K\exp\!\left(2z-\frac\pi2(e^{2|x|}+e^{2|x+s|})\right)\\
 &\le K e^{2z-\pi e^z}\le K e^{2s-\pi e^s}.
\end{align*}
The last step uses $2-\pi e^z<0$ for $z\ge0$. The norm of each weighted translation is bounded by the supremum of its multiplier, so its two-shift contribution at $s_n$ has norm at most $2Ka_nn^2e^{-\pi n}$. Since $\Lambda(n)\le\log n$, the norm series converges. Adjointness interchanges the two shifts, and reflection also interchanges them. This proves self-adjointness and even invariance.

Writing $u=\beta F$, the $n$th quadratic term equals the corresponding term in~\eqref{eq:physicalprime}. Its absolute value is bounded by the summable operator-norm bound times $\norm{F}_2^2=\norm{u}_\omega^2$. This proves both the identification and absolute convergence. The tail estimate follows from~\eqref{eq:rN}; Cauchy--Schwarz and translation invariance of ordinary $L^2$ give the finite-head estimate.
\end{proof}

\begin{remark}
Lemma~\ref{lem:prime} does not assert that $K_P$ is compact. Multiplication by a decaying function followed by a translation need not be compact on a nonatomic $L^2$ space. Nor does~\eqref{eq:rN} assert a uniformly vanishing tail for the unweighted incoming prime rate as the physical point escapes to infinity. The norm, the weight, and the coordinates in the lemma are essential.
\end{remark}

\section{The exact numerator domain}\label{sec:domain}

We first analyze $X$ independently of the bridge. Because $W$ is positive and bounded,
\begin{equation}\label{eq:L2bound}
 \norm{u}_2^2\le\norm{W}_\infty\norm{u}_\omega^2.
\end{equation}
The Gamma form is closed on ordinary $L^2$, for example by the same difference-and-Fatou argument used for $\E$.

\begin{lemma}\label{lem:Xcore}
The space $X$ of~\eqref{eq:X} is Hilbert, and even $C_c^\infty(\R)$ is dense in $X$.
\end{lemma}
\begin{proof}
A sequence Cauchy in $X$ has a weighted $L^2$ limit and an ordinary Gamma-form limit. Equation~\eqref{eq:L2bound} identifies their ordinary $L^2$ limits, proving completeness.

Use the cutoffs $\chi_R$ from Proposition~\ref{prop:core}. For $h_R=(1-\chi_R)u$ the product identity gives
\begin{align*}
 |h_R(x+s)-h_R(x)|^2
 &\le2(1-\chi_R(x+s))^2|u(x+s)-u(x)|^2\\
 &\quad+2|u(x)|^2|\chi_R(x+s)-\chi_R(x)|^2.
\end{align*}
The first term tends to zero in the Gamma integral by dominated convergence. The second is dominated by
$2|u(x)|^2\min\{1,L^2s^2\}$, which is integrable by~\eqref{eq:rhointegrable}. Weighted convergence is another application of dominated convergence. Thus compactly supported vectors are dense in $X$.

For a fixed compactly supported vector, even smooth mollification converges in ordinary $L^2$ and in $\tg$, by the translation argument in Proposition~\ref{prop:core}. The supports remain in one compact interval, where $1/W$ is bounded. Ordinary $L^2$ convergence consequently gives weighted convergence as well.
\end{proof}

For original profiles and physical numerators the coordinates are
\begin{equation}\label{eq:coordinates}
 u=\kappa f,\qquad v=2cf=\frac uW,\qquad
 F=\frac u{\sqrt W}=\sqrt{2c\kappa}\,f.
\end{equation}
Only $f\mapsto F$ is the displayed unitary map from $\He$ to ordinary $\Le$. The map $f\mapsto u$ is an isometry to the weighted space with norm $\norm{\cdot}_\omega$, since
\begin{equation}\label{eq:weightedidentity}
 \norm{\kappa f}_\omega=\norm{f}_\mu.
\end{equation}

\begin{proposition}[Extension of the compact bridge]\label{prop:domain}
Assume \textup{(H)}. Multiplication by $\kappa$ gives a bijection $\D\to X$, with equivalent graph norms, and
\begin{equation}\label{eq:extendedbridge}
 q_0(\kappa f,\kappa g)=\E(f,g)-\tfrac12\ip{f}{g}_\mu,
 \qquad f,g\in\D.
\end{equation}
\end{proposition}
\begin{proof}
Put $C_0=|\cg|\norm{W}_\infty+\norm{K_P}$. On the compact smooth core, hypothesis~(H), Lemma~\ref{lem:prime}, and~\eqref{eq:weightedidentity} imply
\begin{align}
 \tg[\kappa f]&\le\E[f]+(C_0+\tfrac12)\norm{f}_\mu^2,\label{eq:graph1}\\
 \E[f]&\le\tg[\kappa f]+(C_0+\tfrac12)\norm{f}_\mu^2.\label{eq:graph2}
\end{align}
Taking compact-profile approximants from Proposition~\ref{prop:core} and using completeness of $X$ proves that $\kappa\D\subset X$.

Conversely, for $u\in X$ choose even $u_j\in C_c^\infty$ converging in $X$, by Lemma~\ref{lem:Xcore}. Since $\kappa$ is smooth and strictly positive, $f_j=u_j/\kappa$ are legitimate even compact smooth profiles. Equation~\eqref{eq:graph2} makes $(f_j)$ Cauchy in the original energy graph norm. Its weighted limit is $f=u/\kappa$, so closedness of $\E$ gives $f\in\D$. This proves surjectivity and graph-norm equivalence. The Gamma form and the prime form are continuous in $X$, so the quadratic identity passes to the limit; polarization gives~\eqref{eq:extendedbridge}.
\end{proof}

Define $m(u)=\int_\R c(x)u(x)\,dx$. The integral exists absolutely for finite weighted norm, because
\begin{equation}\label{eq:meanbound}
 |m(u)|\le\left(\int_\R c^2W\,dx\right)^{1/2}\norm{u}_\omega
       =\tfrac12\norm{u}_\omega.
\end{equation}
Here $\int c^2W=\tfrac12\int c\kappa=1/4$. Thus $\mu(f)=2m(\kappa f)$, and under (H) the full pole form obeys
\begin{equation}\label{eq:polebridge}
 q[u]:=q_0[u]+2|m(u)|^2
       =\E[f]-\tfrac12\Var_\mu(f),\qquad u=\kappa f\in X.
\end{equation}
The coefficient $2$ in the pole is fixed by normalization. It is not an additional positive contribution available after centering.

\section{The closed pencil and compact resolvent}\label{sec:pencil}

\begin{proposition}[Closed ordinary-$L^2$ realization]\label{prop:closedpencil}
For every $\delta>0$, the form $\ell_\delta$ of~\eqref{eq:ell} is closed and lower bounded on $\Le$. Its self-adjoint representative is characterized by
\begin{equation}\label{eq:Ldef}
 \begin{split}
 u\in\dom(L_\delta),\ L_\delta u=h\quad\Longleftrightarrow\quad
 u\in X,\ h\in\Le,\quad
 \ell_\delta(u,z)=\ip{h}{z}_2\ \text{for all }z\in X.
 \end{split}
\end{equation}
No use of \textup{(H)} is needed.
\end{proposition}
\begin{proof}
Choose an integer $N$ for which $r_N\le\delta/2$ and set
\[
 C_N=|\cg|+2\sum_{2\le n\le N}a_n.
\]
The complete head--tail estimate gives
\begin{equation}\label{eq:coercive}
 \ell_\delta[u]\ge\tg[u]+\frac\delta2\norm{u}_\omega^2-C_N\norm{u}_2^2.
\end{equation}
Its upper absolute bound is a finite multiple of $\norm{u}_X^2$, by Lemma~\ref{lem:prime} and~\eqref{eq:L2bound}. Therefore
$\ell_\delta[u]+(C_N+1)\norm{u}_2^2$ defines a positive norm squared equivalent to $\norm{u}_X^2$. Completeness and density follow from Lemma~\ref{lem:Xcore}. The closed-form representation theorem gives the self-adjoint operator and~\eqref{eq:Ldef}.
\end{proof}

The choice of $N$ in this proof is existential and analytic. The form remains the complete form for every such choice. In distributional notation the operator is
\begin{equation}\label{eq:formalL}
 L_\delta=D_\Gamma-\cg-P+\frac\delta W,
 \qquad
 Pu=\sum_{n\ge2}a_n\bigl(u(\cdot+s_n)+u(\cdot-s_n)\bigr),
\end{equation}
where $D_\Gamma$ is represented by $\tg$. For $u\in X$, the prime row converges in $L^2$ on every fixed compact interval: each shifted local norm is bounded by the local supremum of $\sqrt W$ times $\norm{u}_\omega$, and the theta envelope makes the resulting tails summable. Also $u/W$ is locally $L^2$. The Gamma term is defined distributionally by its form. Operator-domain membership requires that the \emph{sum} in~\eqref{eq:formalL} be represented by an ordinary $L^2$ vector; it does not assert that each term is globally $L^2$. Compact smooth functions are a form core. No operator-core claim for this ordinary-$L^2$ realization is being made.

\begin{lemma}[Logarithmic Fourier coercivity]\label{lem:fourier}
For the unitary Fourier transform,
\begin{equation}\label{eq:multiplier}
 \tg[u]=\int_\R d_\Gamma(\xi)|\widehat u(\xi)|^2\,d\xi,
 \quad
 d_\Gamma(\xi)=2\int_0^\infty\rho(s)(1-\cos(\xi s))\,ds
 \ge\frac12\log(1+4\xi^2).
\end{equation}
\end{lemma}
\begin{proof}
Plancherel's theorem and Tonelli's theorem give the multiplier identity, allowing infinite values. The inequality $1-e^{-2s}\le2s$ gives
\[
 d_\Gamma(\xi)\ge\int_0^\infty e^{-s/2}\frac{1-\cos(\xi s)}s\,ds.
\]
For general $a>0$, differentiating this last integral with $e^{-as}$ in place of $e^{-s/2}$ yields $\int_0^\infty e^{-as}\sin(\xi s)\,ds=\xi/(a^2+\xi^2)$. Integrating back from $\xi=0$ gives $\tfrac12\log(1+\xi^2/a^2)$. Taking $a=1/2$ proves the bound.
\end{proof}

\begin{proposition}[Compactness]\label{prop:compact}
The embedding $X\hookrightarrow\Le$ is compact. Consequently $L_\delta$ has compact resolvent for every $\delta>0$.
\end{proposition}
\begin{proof}
The unit ball of $X$ is uniformly spatially tight, since
\begin{equation}\label{eq:spatialtight}
 \int_{|x|>R}|u(x)|^2\,dx
 \le\sup_{|x|>R}W(x)\,\norm{u}_\omega^2\longrightarrow0
\end{equation}
uniformly on that ball. Lemma~\ref{lem:fourier} gives uniform Fourier tightness:
\begin{equation}\label{eq:fouriertight}
 \int_{|\xi|>M}|\widehat u(\xi)|^2\,d\xi
 \le\frac{\tg[u]}{\tfrac12\log(1+4M^2)}\longrightarrow0.
\end{equation}
Let $P_M$ be the Fourier projection onto $[-M,M]$. The operator
$\one_{[-R,R]}P_M$ is Hilbert--Schmidt: its space--frequency kernel has constant modulus $(2\pi)^{-1/2}$ on a finite rectangle and vanishes elsewhere. Moreover
\[
 \norm{u-\one_{[-R,R]}P_Mu}_2
 \le\norm{\one_{\{|x|>R\}}u}_2+\norm{(I-P_M)u}_2.
\]
Thus the inclusion is a uniform limit on the $X$ unit ball of compact operators. The cutoffs and Fourier projection preserve evenness.

By~\eqref{eq:coercive}, a resolvent $(L_\delta+aI)^{-1}$ with $a>C_N$ maps ordinary $L^2$ boundedly into $X$. Composing this map with the compact inclusion proves compactness of the resolvent. Hence the eigenvalues are real, have finite multiplicity, are bounded below, and tend to $+\infty$.
\end{proof}

This compactness concerns $L_\delta$ on ordinary $L^2$, not the original operator $A_E$ on $L^2(\mu)$. The two Hilbert norms are different.

\section{Actual eigenstates and the negative index}\label{sec:spectrum}

In this section and the rest of the paper, hypothesis~\textup{(H)} is assumed. Proposition~\ref{prop:domain} gives, for $\lambda=1/2-\delta$,
\begin{equation}\label{eq:congruence}
 \ell_\delta(\kappa f,\kappa g)
       =\E(f,g)-\lambda\ip{f}{g}_\mu,
       \qquad f,g\in\D.
\end{equation}
This is a form congruence through a bijection of form domains, not a unitary equivalence of operators on their displayed Hilbert spaces.

\begin{proposition}[Fixed-distance spectral finiteness]\label{prop:finite}
For every $t<1/2$, the spectral projection $\one_{(-\infty,t]}(A)$ has finite rank. The centered restriction satisfies $A_E\ge\varepsilon_0 I$ for some $\varepsilon_0>0$ and
\begin{equation}\label{eq:essential}
 \inf\sigma_{\ess}(A_E)\ge\frac12,
\end{equation}
where the infimum of the empty essential spectrum is $+\infty$.
\end{proposition}
\begin{proof}
For $t<0$ the assertion follows from $A\ge0$. Let $0\le t<1/2$, and take $\delta=1/2-t$. Every $f$ in the range of $\one_{[0,t]}(A)$ lies in $\dom(A)$ and satisfies $\E[f]\le t\norm{f}_\mu^2$. Equation~\eqref{eq:congruence} sends this range injectively into a subspace of $X$ on which $\ell_\delta\le0$.

A lower-bounded compact-resolvent operator has only finitely many nonpositive eigenvalues. Its nonpositive spectral subspace is therefore finite dimensional. Any subspace of its form domain on which the form is nonpositive has dimension at most that finite number: otherwise a nonzero vector in the subspace would be orthogonal to the entire nonpositive spectral subspace, where the form is strictly positive. Applying this observation to $L_\delta$ proves finite rank.

The centered restriction has no zero eigenvector, by Proposition~\ref{prop:core}. Apply the finite-rank statement at any one $t_0\in(0,1/2)$, for example $t_0=1/4$. On its centered finite-dimensional spectral range, if nonzero, the smallest eigenvalue is strictly positive; on the complementary range the operator is bounded below by $t_0$. This proves a positive lower bound. Finally, finite rank below every $t<1/2$ excludes essential spectrum below $1/2$ and permits accumulation of subthreshold eigenvalues only at $1/2$.
\end{proof}

The proposition is a consequence of the compact pencil and the bridge. It does not require an additional estimate for primes in short intervals, nor does it assert equality in~\eqref{eq:essential}.

\begin{proposition}[Kernel correspondence]\label{prop:kernel}
For $0<\delta<1/2$ and $\lambda=1/2-\delta$, equation~\eqref{eq:mainkernel} holds as an exact linear correspondence, including multiplicities.
\end{proposition}
\begin{proof}
If $u\in\ker L_\delta$, then $u\in X$ and $f=u/\kappa\in\D$. Testing~\eqref{eq:Ldef} against every $z=\kappa g\in X$ and using~\eqref{eq:congruence} gives
\[
 \E(f,g)=\lambda\ip{f}{g}_\mu\qquad(g\in\D).
\]
The representation criterion~\eqref{eq:Adef} proves $f\in\dom(A)$ and $Af=\lambda f$. Since $\lambda>0$ and $A1=0$, self-adjointness implies $\mu(f)=0$. Thus $f\in\dom(A_E)$ is a genuine centered eigenstate.

Conversely, for such an eigenstate $f$, the vector $u=\kappa f$ belongs to $X$ and~\eqref{eq:congruence} vanishes against every $z\in X$. Its ordinary-$L^2$ representing vector is zero. Hence $u\in\dom(L_\delta)$ and $L_\delta u=0$.
\end{proof}

No pointwise eigen-row or additional regularity of an eigenstate is used here. The correspondence identifies kernels at the indicated parameter; it does not identify nonzero eigenvalues of $L_\delta$ with eigenvalues of $A_E$.

\begin{proposition}[Exact form index]\label{prop:index}
For $0<\delta<1/2$ and $\lambda=1/2-\delta$,
\begin{align}
 n_-(L_\delta)&=1+\rank\one_{(0,\lambda)}(A_E),\label{eq:index}\\
 \dim\ker L_\delta&=\dim\ker(A_E-\lambda).\label{eq:nullity}
\end{align}
For $u\in X$, put $a=2m(u)$ and $u_0=u-a\kappa$. Then
\begin{align}
 m(u_0)&=0,\qquad
 \norm{u}_\omega^2=|a|^2+\norm{u_0}_\omega^2,\label{eq:decompnorm}\\
 \ell_\delta[u]&=-\lambda|a|^2+
              \E[u_0/\kappa]-\lambda\norm{u_0}_\omega^2.\label{eq:decompform}
\end{align}
\end{proposition}
\begin{proof}
Write $f=u/\kappa=a+g$, where $a=\mu(f)$ and $\mu(g)=0$. Constants have zero energy, and polarization gives $\E(1,g)=0$. Also $\norm{1}_\mu=1$. The orthogonal decomposition in $L^2(\mu)$ and~\eqref{eq:congruence} prove~\eqref{eq:decompnorm}--\eqref{eq:decompform}.

A bijective linear change of form domain preserves the maximal dimension of a strictly negative form subspace. On the constant line the right side of~\eqref{eq:decompform} is strictly negative. On the centered subspace its negative index is the rank of the spectral projection of $A_E$ below $\lambda$, by the spectral theorem. Proposition~\ref{prop:finite} and the positive gap make this the finite rank in~\eqref{eq:index}. For a compact-resolvent operator its form negative index equals its number of negative eigenvalues. The nullity identity is Proposition~\ref{prop:kernel}.
\end{proof}

\begin{corollary}[Failure of the half bound]\label{cor:failure}
The global inequality $\E[f]\ge\tfrac12\Var_\mu(f)$ on $\D$ fails if and only if there exists an actual eigenpair
\[
 A_Ef=\lambda f,\qquad f\in\dom(A_E)\setminus\{0\},\qquad0<\lambda<1/2.
\]
\end{corollary}
\begin{proof}
Subtract the mean of a violating profile to obtain $g\ne0$ with $\E[g]/\norm{g}_\mu^2<1/2$. Choose $t$ strictly between this quotient and $1/2$. The spectral projection $\one_{(-\infty,t]}(A_E)$ is nonzero, since otherwise the spectral theorem would give $\E[g]\ge t\norm{g}_\mu^2$. It is finite rank by Proposition~\ref{prop:finite}, so its invariant range contains an eigenvector. Its eigenvalue is positive by the gap and at most $t<1/2$. The converse follows by evaluating the energy on the eigenvector.
\end{proof}

\section{The mean and the operator-domain obstruction}\label{sec:mean}

The decomposition in Proposition~\ref{prop:index} is orthogonal in the weighted norm, but not in ordinary physical $L^2$. Two concrete facts make the distinction unavoidable.

\begin{proposition}[The mean is unbounded in ordinary $L^2$]\label{prop:mean}
The functional $m$ is continuous on $X$ and on the weighted space, with bound~\eqref{eq:meanbound}, but is not bounded in the ordinary $L^2$ norm, even on even $C_c^\infty(\R)$. Its zero-mean subspace in $X$ is closed in $X$ and dense in ordinary $\Le$.
\end{proposition}
\begin{proof}
Only unboundedness and density remain to be shown. Fix a nonzero, nonnegative, even $\phi\in C_c^\infty(\R)$ and put $M=m(\phi)>0$. For sufficiently large $R$, the two translates in
\[
 \psi_R(x)=\phi(x-R)+\phi(x+R),\qquad
 b_R(x)=\frac{\psi_R(x)}{2\cosh(R/2)M}
\]
have disjoint supports. The addition formula for the hyperbolic cosine and evenness of $\phi$ give $m(\psi_R)=2\cosh(R/2)M$. Hence
\begin{equation}\label{eq:unboundedmean}
 m(b_R)=1,\qquad
 \norm{b_R}_2^2=\frac{\norm{\phi}_2^2}{2\cosh^2(R/2)M^2}\longrightarrow0.
\end{equation}
Every $b_R$ is in the compact smooth numerator core. Thus neither $m$ nor the quadratic pole $2|m|^2$ has an ordinary-$L^2$ bound. For any even compact smooth $u$, the corrected function $u-m(u)b_R$ has zero mean and converges to $u$ in ordinary $L^2$. Density of the compact smooth core in $\Le$ proves the last claim. No smallness of $\norm{b_R}_X$ or $\norm{b_R}_\omega$ is asserted.
\end{proof}

\begin{proposition}[The theta kernel is generally not in the operator domain]\label{prop:kdomain}
For every $\delta>0$, $\kappa\in X$. If $\delta\ne1/2$, then $\kappa\notin\dom(L_\delta)$. At $\delta=1/2$,
\begin{equation}\label{eq:halfkernel}
 L_{1/2}\ge0,\qquad \ker L_{1/2}=\C\kappa.
\end{equation}
For $0<\delta<1/2$, $\kappa$ is nevertheless a strictly negative form direction.
\end{proposition}
\begin{proof}
Since $1\in\D$, Proposition~\ref{prop:domain} admits $\kappa$ to $X$ and gives $\norm{\kappa}_\omega=1$, $m(\kappa)=1/2$. Put $\lambda=1/2-\delta$, now allowing either sign. Equation~\eqref{eq:congruence} gives the exact polarized identity
\begin{equation}\label{eq:sourcerow}
 \ell_\delta(\kappa,z)=-2\lambda\overline{m(z)},\qquad z\in X.
\end{equation}
If $\kappa\in\dom(L_\delta)$, then compact smooth even tests in~\eqref{eq:Ldef} force its ordinary-$L^2$ representative to equal $-2\lambda c$ as an even distribution. For $\lambda\ne0$ this is impossible, since $c\notin L^2(\R)$. When $\lambda=0$, the form is $\E[u/\kappa]\ge0$ and its zero vectors are precisely multiples of $\kappa$, by Proposition~\ref{prop:core}. Finally $\ell_\delta[\kappa]=-\lambda<0$ for $0<\delta<1/2$.
\end{proof}

Thus a negative \emph{form direction} is not necessarily an operator-domain vector, much less an ordinary-$L^2$ ground eigenvector. The mean-zero space is not an ordinary-$L^2$ closed codimension-one hyperplane. Moreover $\int u\overline\kappa\,dx=0$ is different from $m(u)=0$: the two functionals on compact tests would be proportional only if $\kappa$ were a scalar multiple of $c$, contrary to their decay. Removing $\kappa$ by an ordinary-$L^2$ orthogonal projection therefore does not implement centering of the original problem.

\section{The second eigenvalue and the remaining sign}\label{sec:sign}

The compact-resolvent construction gives a precise family of sign questions, but it supplies no positive sign for the second eigenvalue throughout the required parameter interval.

\begin{proposition}[Equivalent sign conditions]\label{prop:sign}
Under \textup{(H)}, the following statements are equivalent:
\begin{enumerate}
\item $\E[f]\ge\tfrac12\Var_\mu(f)$ for every $f\in\D$.
\item For every $0<\delta<1/2$, $\eta_2(\delta)>0$.
\item For every $0<\delta<1/2$ and every nonzero $u\in X$ with $m(u)=0$, one has $\ell_\delta[u]>0$.
\item $\ker L_\delta=\{0\}$ for every $0<\delta<1/2$.
\end{enumerate}
\end{proposition}
\begin{proof}
If the first statement holds, then on centered vectors~\eqref{eq:congruence} gives $\ell_\delta[u]\ge\delta\norm{u}_\omega^2>0$. It also gives no centered spectrum of $A_E$ at or below $\lambda=1/2-\delta$. Equations~\eqref{eq:index}--\eqref{eq:nullity} then show that $L_\delta$ has exactly one negative eigenvalue and no zero eigenvalue; compact resolvent gives $\eta_2(\delta)>0$.

Conversely, if $\eta_2(\delta)>0$, the index and nullity identities exclude centered spectrum at or below $\lambda$, proving strict positivity on the centered form subspace. If the third statement holds for all $\delta$, letting $\delta\downarrow0$ for each fixed centered $u$ gives $q_0[u]\ge0$. By~\eqref{eq:polebridge} and subtraction of the original mean, this is the first statement. Finally, the equivalence of the first and fourth statements follows from Corollary~\ref{cor:failure} and Proposition~\ref{prop:kernel}.
\end{proof}

Two valid estimates illustrate the limits of the construction. First, min--max on the common domain $X$ and~\eqref{eq:L2bound} give
\begin{equation}\label{eq:monotonicity}
 \eta_j(\delta_2)\ge\eta_j(\delta_1)
       +\frac{\delta_2-\delta_1}{\norm{W}_\infty},
       \qquad\delta_2>\delta_1>0.
\end{equation}
At $\delta=1/2$, Proposition~\ref{prop:kdomain} and compact resolvent give $\eta_2(1/2)>0$. However, decreasing $\delta$ decreases the form. Inequality~\eqref{eq:monotonicity} does not prohibit a second eigenvalue from crossing zero at a smaller positive parameter.

Second, suppose $u\in\ker L_\delta$ with $0<\delta<1/2$ and normalize $\norm{u}_\omega=1$. Then
\begin{equation}\label{eq:stateidentity}
 \tg[u]+\delta=\cg\norm{u}_2^2+p[u].
\end{equation}
Choose $N$ with $r_N\le\delta/2$ as in Proposition~\ref{prop:closedpencil}. The same complete-tail estimate yields
\begin{equation}\label{eq:statebound}
 \tg[u]+\frac\delta2\le C_N\norm{u}_2^2,
 \qquad \norm{u}_2^2\ge\frac{\delta}{2C_N}.
\end{equation}
Thus a normalized fixed-gap candidate cannot lose all its ordinary physical $L^2$ mass. The estimate is not contradictory: reducing the tail by enlarging $N$ also enlarges the head constant $C_N$. No estimate reversing this balance is proved. Testing the kernel equation against $\kappa$ gives only $m(u)=0$, by~\eqref{eq:sourcerow}; this recovers centering and provides no second independent sign constraint.

\section{Scope and conclusion}

The complete weighted prime estimate, the closed ordinary-$L^2$ pencil, and its compact resolvent are established directly from the theta envelope and the Gamma form. Under the separately stated bridge~\textup{(H)}, the original and physical form domains coincide through multiplication by $\kappa$, and the correspondence reaches actual operator-domain eigenstates. Negative-index accounting then identifies the exact second-eigenvalue sign needed for the global variance half bound.

Several boundaries remain important. The correspondence is not a unitary identification of the full spectra. Compactness of the pencil does not imply compactness of the weighted prime operator or compact resolvent of the original energy generator. The essential spectral bottom of $A_E$ is bounded below by $1/2$ here; equality or an infinite-multiplicity half-eigenspace would require additional hypotheses not assumed in this paper. The compact bridge itself is not reproved. Finally, no subthreshold eigenstate is excluded by the bounds above, and the global half bound and the Riemann hypothesis remain unproved.

\begin{thebibliography}{9}
\bibitem{DLMF}
F.~W.~J. Olver, A.~B. Olde Daalhuis, D.~W. Lozier, B.~I. Schneider,
R.~F. Boisvert, C.~W. Clark, B.~R. Miller, B.~V. Saunders, H.~S. Cohl,
and M.~A. McClain, eds.,
\emph{NIST Digital Library of Mathematical Functions},
Section~20.7, Eq.~20.7.32, Jacobi transformation.
\url{https://dlmf.nist.gov/20.7.E32}.

\bibitem{Kato}
T. Kato, \emph{Perturbation Theory for Linear Operators},
second edition, Classics in Mathematics, Springer, Berlin, 1995.
Chapter~VI, sesquilinear forms and associated operators.
\url{https://doi.org/10.1007/978-3-642-66282-9}.
\end{thebibliography}
\end{document}
